% MACROS:
% \newcommand{\dn}{\mathord{\downarrow}}
% \newcommand{\upp}{\mathord{\uparrow}}
% \newcommand{\LMN}{\mathrm{LMN}}
% \newcommand{\UMN}{\mathrm{UMN}}
% \newcommand{\pos}{\mathrm{pos}}
% \newcommand{\nodeof}{\mathrm{node}}
% \newcommand{\pk}{\mathrm{pk}}

\section{From Pairs at Distance Two to a Zigzag Problem}
\label{sec:zz-zigzag}

An \emph{instance of the open case} is a triple $(G, u, v)$ with $G \in \DC$ and vertices $u, v$ of $G$ with $d(u, v) = 2$ and $c(u, v) \geq 3$.
By \cref{prop:easy-pairs}(iv), we have $\cp(u, v) \in \{c(u, v), c(u, v) + 1\}$, so the question is whether $\cp(u, v) = c(u, v)$.
\Cref{thm:zz-reduction} maps such an instance to an instance of a problem on a preordered set with the same answer, and \cref{lem:zz-answers} rewrites that problem in the form solved in \cref{sec:pa-polynomial}.
In this section and in \cref{sec:pa-polynomial}, $m$ counts the switches of an instance of \cref{def:zz-general}, and the number of edges of $G$ is written $|E|$.

\subsection{Preorders and the Problem ZZ}
\label{sec:zz-preorders}

\begin{definition}[Preorders]
\label{def:zz-preorder}
A \emph{preorder} $\leq$ on a finite set is a reflexive and transitive relation.
We write $x < y$ if $x \leq y$ and not $y \leq x$.
We write $x \sim y$, and call $x$ and $y$ \emph{tied}, if $x \leq y$ and $y \leq x$, and $[x]$ is the class of $x$ under this equivalence relation.
We write $x \parallel y$, and call $x$ and $y$ \emph{incomparable}, if neither $x \leq y$ nor $y \leq x$.
For a set $S$, we put $\dn S = \{x : x \leq y \text{ for some } y \in S\}$, $\upp S = \{x : x \geq y \text{ for some } y \in S\}$, $\dn x = \dn \{x\}$ and $\upp x = \upp \{x\}$.
Hence $\dn x \setminus [x] = \{y : y < x\}$ and $\upp x \setminus [x] = \{y : y > x\}$.
A \emph{down-set} $D$ satisfies $\dn D = D$, and an \emph{up-set} $D$ satisfies $\upp D = D$.
Complements of down-sets are up-sets, intersections and unions of down-sets are down-sets, $\dn S$ is a down-set, and the complement of $\upp S$ is a down-set.
If $x \leq y < z$ or $x < y \leq z$, then $x < z$, because $z \leq x$ would give $z \leq y$ in the first case and $y \leq x$ in the second.
For every two elements $x, y$, exactly one of $x \leq y$, $y < x$ and $x \parallel y$ holds.
\end{definition}

\begin{definition}[ZZ and P2]
\label{def:zz-zz}
An instance of \emph{ZZ} is a tuple $(Q, \leq, R, u, v)$ that consists of a finite set $Q$ with a preorder $\leq$ such that $u \leq x \leq v$ for every $x \in Q$, and a set $R$ of pairs $\{a, b\} \subseteq Q$ with $a < b$ such that $\{u, v\} \notin R$.
The pairs of $R$ are \emph{real edges}, and we write $ab$ for a real edge with $a < b$.
For a down-set $D$ of $Q$ with $u \in D$ and $v \notin D$, the graph $G_D$ on $Q$ has as its edges the real edges with exactly one end in $D$.
The instance is a \emph{yes-instance} if $u$ and $v$ are connected in $G_D$ for some such $D$.
It satisfies \emph{P2} if every two incomparable elements $x, y$ have a common $R$-neighbor $z$ with $z < x$ and $z < y$ (a \emph{lower witness}) or with $z > x$ and $z > y$ (an \emph{upper witness}).
\end{definition}

A path of $G_D$ alternates between $D$ and $Q \setminus D$, which gives ZZ its name (zigzag).

\subsection{The General Switch Form}
\label{sec:zz-switch-form}

\begin{definition}[General switch form]
\label{def:zz-general}
An instance has a finite set $\Omega$ of \emph{nodes} with two \emph{terminals} $s, t \in \Omega$, a finite set $M$ of \emph{switches} with a preorder $\leq$, and a set of middle edges.
We write $m = |M|$.
Every switch $w$ has an \emph{E-side} and an \emph{F-side}, with node sets $E_w \subseteq \Omega$ and $F_w \subseteq \Omega$.
Sides are labeled, that is, a side is a switch together with E or F, and we denote a side by its node set, $E_w$ or $F_w$.
A \emph{middle edge} is an ordered pair $(z, w)$ of switches with $z < w$, and it has a \emph{private node} $\rho(z, w) \in \Omega$ that lies in $F_z$ and in $E_w$ and in no other side.
The sets $\LMN(w) = \{z : (z, w) \text{ is a middle edge}\}$ and $\UMN(w) = \{z : (w, z) \text{ is a middle edge}\}$ consist of the lower and the upper middle neighbors of $w$.
The instance satisfies \emph{P2} if every two incomparable switches $x, y$ have a common E-node ($E_x \cap E_y \neq \emptyset$), or a common F-node ($F_x \cap F_y \neq \emptyset$), or a \emph{lower middle witness} $z \in \LMN(x) \cap \LMN(y)$, or an \emph{upper middle witness} $z \in \UMN(x) \cap \UMN(y)$.
\end{definition}

\begin{definition}[Mode vectors, conflict pairs, valid paths]
\label{def:zz-valid}
A \emph{mode vector} is a down-set $D$ of $(M, \leq)$.
Switches in $D$ are in \emph{F-mode}, the others in \emph{E-mode}.
The mode vector $D$ \emph{picks} $F_w$ for $w \in D$ and $E_w$ for $w \notin D$.
It \emph{connects} $s$ and $t$ if $s = t$ or there are picked sides $X_1, \dots, X_h$ with $h \geq 1$, $s \in X_1$, $t \in X_h$ and $X_i \cap X_{i+1} \neq \emptyset$ for $i < h$.
The instance is a \emph{yes-instance} if some mode vector connects $s$ and $t$.
A \emph{conflict pair} is a pair of sides $\{E_w, F_z\}$ with $w \leq z$, in particular $\{E_w, F_w\}$.
A \emph{valid path} is a sequence $P = X_1, \dots, X_q$ with $q \geq 1$ of sides of pairwise distinct switches, with $s \in X_1$, $t \in X_q$ and $X_i \cap X_{i+1} \neq \emptyset$ for $i < q$, no two of which form a conflict pair.
Its sides are its \emph{members}.
A switch $w$ is \emph{on} $P$ if one of its sides is a member, and then $\pos(w)$ is the index of that member.
We write $W_E$ and $W_F$ for the sets of switches whose E-side and whose F-side, respectively, is a member, and we put $U_P = \upp W_E$ and $L_P = \dn W_F$.
A valid path is \emph{shortest} if no valid path has fewer members.
\end{definition}

\subsection{Reduction of the Open Case to ZZ with P2}
\label{sec:zz-reduction}

Throughout this subsection, $(G, u, v)$ is an instance of the open case with $n$ vertices.

\begin{lemma}[Bridges]
\label{lem:zz-bridges}
Let $K$ be the component containing $u$ of the graph obtained from $G$ by removing all bridges.
(a) We have $v \in V(K)$.
(b) For $x, y \in V(K)$, the $x$--$y$ paths of $G$ are exactly the $x$--$y$ paths of $K$.
Hence $d_K(x, y) = d(x, y)$, $c_K(x, y) = c(x, y)$, and an edge set is a minimum $x$--$y$ cut of $G$ if and only if it is a minimum $x$--$y$ cut of $K$.
(c) Some minimum $u$--$v$ cut of $G$ contains a $u$--$v$ path if and only if some minimum $u$--$v$ cut of $K$ does.
(d) The graph $K$ is connected, has no bridge and lies in $\DC$, and $d_K(u, v) = 2$, $c_K(u, v) \geq 3$.
\end{lemma}

\begin{proof}
A bridge $b = pp'$ of the connected graph $G$ lies on no cycle, and $G \setminus \{b\}$ has exactly two components $Y_p \ni p$ and $Y_{p'} \ni p'$.
If a path with ends $x$ and $y$ uses $b$, traversed from $x$ to $p$, then to $p'$ and then to $y$, then its two sub-paths avoid $b$, so $x \in Y_p$, $y \in Y_{p'}$, and $b$ separates $x$ from $y$.
Call this fact $(*)$.

(a) A $u$--$v$ path through a bridge $b$ would make $\{b\}$ a $u$--$v$ cut by $(*)$, against $c(u, v) \geq 3$.
As $G$ is connected, a $u$--$v$ path exists.
It uses no bridge, so it lies in the component $K$ of $u$.

(b) Let $x, y \in V(K)$.
They are joined by a path of $K$, which avoids every bridge $b$, so $x$ and $y$ lie in one component of $G \setminus \{b\}$, and by $(*)$ no $x$--$y$ path of $G$ uses $b$.
Hence every $x$--$y$ path of $G$ avoids all bridges and, starting in $K$, lies in $K$.
Conversely, every path of $K$ is a path of $G$.
So the distances agree, and an edge set $C \subseteq E$ meets every $x$--$y$ path of $G$ if and only if $C \cap E(K)$ meets every $x$--$y$ path of $K$.
Therefore $c(x, y) = c_K(x, y)$, a minimum cut of $G$ lies in $E(K)$ because otherwise $C \cap E(K)$ is a smaller cut, and a minimum cut of $K$ is a cut of $G$ of the same size.

(c) By (b) with $x = u$ and $y = v$, the minimum cuts and the $u$--$v$ paths are the same in $G$ and in $K$.

(d) The graph $K$ is a component, hence connected.
An edge $e$ of $K$ is not a bridge of $G$, so it lies on a cycle of $G$.
That cycle contains no bridge, so it lies in the component $K$, and $e$ is not a bridge of $K$.
Membership in $\DC$ is a condition on $d$ and $c$ of pairs of vertices, which agree by (b).
So $K \in \DC$, $d_K(u, v) = 2$ and $c_K(u, v) = c(u, v) \geq 3$.
\end{proof}

From now on, $K$ is the graph of \cref{lem:zz-bridges}, and we put
\[
Q = \{u\} \cup \{x \in V(K) : c_K(u, x) \geq 3\}.
\]
We have $v \in Q$ because $c_K(u, v) \geq 3$.
For distinct $x, y \in V(K)$ and a set $S \subseteq V(K)$ with $x \in S$ and $y \notin S$, the set $\delta_K(S)$ of edges of $K$ with exactly one end in $S$ is an $x$--$y$ cut of $K$, and by \cref{lem:minimum-cuts} applied to $K$, every minimum $x$--$y$ cut of $K$ has this form.
We also call a set $S \subseteq V(K)$ with $u \in S$, $v \notin S$ and $|\delta_K(S)| = c_K(u, v)$ a \emph{minimum $u$--$v$ cut of $K$}, and we say that it \emph{contains} a path $\pi$ if $\pi \subseteq \delta_K(S)$.

\begin{lemma}[Paths inside minimum cuts lie in $Q$]
\label{lem:zz-paths-in-q}
(a) Let $S$ be a minimum $u$--$v$ cut of $K$ that contains a $u$--$v$ path $\pi$, and let $Y \subseteq V(K) \setminus \{u, v\}$ with $|\delta_K(Y)| \leq 2$.
Then $\pi$ has no vertex in $Y$.
(b) Every vertex of such a path $\pi$ lies in $Q$.
(c) Every two distinct $x, y \in Q$ satisfy $c_K(x, y) \geq 3$ and $d_K(x, y) \leq 2$.
\end{lemma}

\begin{proof}
(a) Suppose that $\pi$ has a vertex in $Y$.
It starts at $u \notin Y$ and ends at $v \notin Y$, so the edge by which it first enters $Y$ and the edge by which it last leaves $Y$ are two distinct edges of $\delta_K(Y)$.
As $|\delta_K(Y)| \leq 2$, we have $\delta_K(Y) = \{e_1, e_2\}$ with $e_i = x_i y_i$, $x_i \in Y$, $y_i \notin Y$, and $e_1, e_2 \in \pi \subseteq \delta_K(S)$.
Let $J$ be the set of edges of $\delta_K(S)$ with no end in $Y$.
Then $|\delta_K(S)| \geq |J| + 2$.
The sets $S \setminus Y$ and $S \cup Y$ contain $u$ and not $v$.
An edge with no end in $Y$ lies in $\delta_K(S \setminus Y)$, and in $\delta_K(S \cup Y)$, if and only if it lies in $\delta_K(S)$, that is, in $J$.
An edge with both ends in $Y$ lies in neither of the two sets.
The edge $e_i$ lies in $\delta_K(S \setminus Y)$ if and only if $y_i \in S$, and in $\delta_K(S \cup Y)$ if and only if $y_i \notin S$.
Hence $|\delta_K(S \setminus Y)| = |J| + |\{i : y_i \in S\}|$ and $|\delta_K(S \cup Y)| = |J| + |\{i : y_i \notin S\}|$.
The two last terms sum to two, so one of the two sets has at most $|J| + 1 < |\delta_K(S)| = c_K(u, v)$ boundary edges.
This set gives a $u$--$v$ cut with fewer than $c_K(u, v)$ edges, a contradiction.

(b) The ends $u$ and $v$ lie in $Q$.
Let $x$ be another vertex of $\pi$, and suppose that $x \notin Q$, so $c_K(u, x) \leq 2$.
By \cref{lem:minimum-cuts}, there is a set $Y$ with $x \in Y$, $u \notin Y$ and $|\delta_K(Y)| = c_K(x, u) \leq 2$.
If $v \in Y$, then $\delta_K(Y)$ is a $u$--$v$ cut with at most $2 < c_K(u, v)$ edges.
So $v \notin Y$ and $Y \subseteq V(K) \setminus \{u, v\}$, and $x \in Y$ contradicts (a).

(c) If $x = u$ or $y = u$, then $c_K(x, y) \geq 3$ by the definition of $Q$.
Otherwise $c_K(x, u) \geq 3$ and $c_K(u, y) \geq 3$ give $c_K(x, y) \geq 3$ by the transitivity in \cref{prop:class}, applied to $K$.
As $K \in \DC$ by \cref{lem:zz-bridges}(d), we have $d_K(x, y) \leq 2$.
\end{proof}

We write $K - Q$ for the graph obtained from $K$ by removing the vertices of $Q$.

\begin{lemma}[Ears]
\label{lem:zz-ears}
Call a set $Y \subseteq V(K)$ an \emph{ear set} if $Y \neq \emptyset$, $Y \cap Q = \emptyset$ and $|\delta_K(Y)| = 2$.
(a) Every vertex of $V(K) \setminus Q$ lies in an ear set.
(b) If two ear sets intersect or are joined by an edge, their union is an ear set.
(c) The inclusion-maximal ear sets, called \emph{ears}, are exactly the vertex sets of the components of $K - Q$.
They partition $V(K) \setminus Q$, no edge joins two ears, and every ear $Y$ has $\delta_K(Y) = \{pa, p'a'\}$ with $a, a' \in Y$ and $p, p' \in Q$, where $a = a'$ or $p = p'$ is allowed.
The vertices $p$ and $p'$ are the \emph{attachments} of $Y$.
(d) Let $E[Y] = E(K[Y]) \cup \delta_K(Y)$ for an ear $Y$.
If the attachments $p, p'$ of $Y$ are distinct and $S \subseteq V(K)$ contains exactly one of them, then $\delta_K(S) \cap E[Y] \neq \emptyset$.
Every edge of $K$ lies either in $K[Q]$ or in $E[Y]$ for exactly one ear $Y$.
\end{lemma}

\begin{proof}
Every set $Z$ with $\emptyset \neq Z \subsetneq V(K)$ has $|\delta_K(Z)| \geq 2$, with at least one edge since $K$ is connected, and a single edge would be a bridge, which $K$ does not have by \cref{lem:zz-bridges}(d).
Call this fact $(**)$.

(a) Let $x \in V(K) \setminus Q$.
Then $x \neq u$ and $c_K(u, x) \leq 2$.
By \cref{lem:minimum-cuts}, there is a set $Y$ with $x \in Y$, $u \notin Y$ and $|\delta_K(Y)| = c_K(x, u) \leq 2$, and $|\delta_K(Y)| = 2$ by $(**)$.
If $x' \in Y \cap Q$, then $x' \neq u$ and $\delta_K(Y)$ separates $x'$ from $u$, so $c_K(u, x') \leq 2$, against $x' \in Q$.
So $Y$ is an ear set containing $x$.

(b) Let $Y$ and $Y'$ be ear sets.
If $Y \cap Y' \neq \emptyset$, then $Y \cup Y'$ and $Y \cap Y'$ are non-empty and miss $u$, so by $(**)$ and \cref{lem:uncrossing} applied to $K$,
\[
4 \leq |\delta_K(Y \cup Y')| + |\delta_K(Y \cap Y')| \leq |\delta_K(Y)| + |\delta_K(Y')| = 4,
\]
whence $|\delta_K(Y \cup Y')| = 2$.
If $Y \cap Y' = \emptyset$ and $j \geq 1$ edges join them, then $|\delta_K(Y \cup Y')| = 4 - 2j \leq 2$, and $|\delta_K(Y \cup Y')| \geq 2$ by $(**)$.
In both cases $Y \cup Y'$ misses $Q$.

(c) By (b), two distinct maximal ear sets are disjoint and not joined by an edge, because their union would be a larger ear set.
By (a), every vertex of $V(K) \setminus Q$ lies in a maximal ear set.
Let $Y$ be an ear.
Distinct components $Z$ of $K[Y]$ have disjoint sets $\delta_K(Z) \subseteq \delta_K(Y)$, because no edge joins $Z$ to $Y \setminus Z$, and each of these sets has size at least two by $(**)$.
As $|\delta_K(Y)| = 2$, the graph $K[Y]$ is connected.
The outer end of an edge of $\delta_K(Y)$ lies outside $Y$ and in no other ear, so it lies in $Q$.
Hence the ears are connected subsets of $K - Q$, pairwise non-adjacent, covering $V(K) \setminus Q$, so they are the vertex sets of the components of $K - Q$.

(d) Let $\delta_K(Y) = \{pa, p'a'\}$ with $p \neq p'$.
The walk that goes from $p$ to $a$, along a path from $a$ to $a'$ in the connected graph $K[Y]$, and from $a'$ to $p'$, is a $p$--$p'$ path with edges in $E[Y]$.
Since $S$ contains exactly one of its ends, one of its edges lies in $\delta_K(S)$.
An edge with an end in an ear $Y$ has its other end in $Y$ or in $Q$ by (c), so it lies in $E[Y]$ and in no $E[Y']$ with $Y' \neq Y$.
An edge with no end in an ear lies in $K[Q]$.
\end{proof}

\begin{lemma}[The multigraph $H$]
\label{lem:zz-multigraph}
Let $H$ be the multigraph on $Q$ whose edges are the edges of $K[Q]$, called \emph{real edges}, and, for every ear with distinct attachments $p$ and $p'$, one edge $pp'$, called its \emph{ear edge}.
The set $\delta_H(B)$ of edges with exactly one end in $B \subseteq Q$ is counted with multiplicity.
A set $B \subseteq Q$ with $u \in B$, $v \notin B$ and least $|\delta_H(B)|$ is a \emph{minimum $u$--$v$ cut of $H$}, and $c_H(u, v)$ is that least value.
For every ear $Y$ fix an order $(p_Y, p'_Y)$ of its attachments, and for $B \subseteq Q$ put $A(B) = B \cup \bigcup \{Y : Y \text{ is an ear with } p_Y \in B\}$.
Then the following hold.
(a) We have $A(B) \cap Q = B$ and $|\delta_K(A(B))| = |\delta_H(B)|$ for every $B \subseteq Q$.
(b) We have $|\delta_K(S)| \geq |\delta_H(S \cap Q)|$ for every $S \subseteq V(K)$.
(c) We have $c_H(u, v) = c_K(u, v)$.
For every minimum $u$--$v$ cut $S$ of $K$, the set $S \cap Q$ is a minimum $u$--$v$ cut of $H$, and for every minimum $u$--$v$ cut $B$ of $H$, the set $A(B)$ is a minimum $u$--$v$ cut of $K$.
(d) Some minimum $u$--$v$ cut of $K$ contains a $u$--$v$ path if and only if some minimum $u$--$v$ cut $B$ of $H$ has a $u$--$v$ path of real edges all of which lie in $\delta_H(B)$.
(e) Every two distinct $x, y \in Q$ are joined by a real edge or by an ear edge, or have a common neighbor $z \in Q$ in $K$, so that $xz$ and $yz$ are real edges.
\end{lemma}

\begin{proof}
By \cref{lem:zz-ears}(d), every edge of $K$ is a real edge or lies in $E[Y]$ for exactly one ear $Y$.

(a) The equality $A(B) \cap Q = B$ holds since ears miss $Q$.
A real edge lies in $\delta_K(A(B))$ if and only if it lies in $\delta_H(B)$.
Edges of $K[Y]$ never lie in $\delta_K(A(B))$, since $Y \subseteq A(B)$ or $Y \cap A(B) = \emptyset$.
Let $\delta_K(Y) = \{p_Y a, p'_Y a'\}$.
If $p_Y \in B$, then $Y \subseteq A(B)$, the edge $p_Y a$ is not cut, and $p'_Y a'$ is cut if and only if $p'_Y \notin B$.
If $p_Y \notin B$, then $Y \cap A(B) = \emptyset$, the edge $p_Y a$ is not cut, and $p'_Y a'$ is cut if and only if $p'_Y \in B$.
So the ear contributes one edge exactly when $B$ contains exactly one of $p_Y$ and $p'_Y$.
This needs $p_Y \neq p'_Y$, and then it is exactly the case that its ear edge lies in $\delta_H(B)$.
Otherwise the ear contributes no edge.

(b) Every real edge of $\delta_H(S \cap Q)$ lies in $\delta_K(S)$.
For every ear edge of $\delta_H(S \cap Q)$, the set $S$ contains exactly one attachment of its ear $Y$, so $\delta_K(S)$ contains an edge of $E[Y]$ by \cref{lem:zz-ears}(d).
These edges are distinct for distinct ears and are not real edges.

(c) For a minimum $u$--$v$ cut $B$ of $H$, the set $A(B)$ contains $u$ and not $v$ by (a), so $c_K(u, v) \leq |\delta_K(A(B))| = |\delta_H(B)| = c_H(u, v)$.
For a minimum $u$--$v$ cut $S$ of $K$, the set $S \cap Q$ contains $u$ and not $v$, so $c_H(u, v) \leq |\delta_H(S \cap Q)| \leq |\delta_K(S)| = c_K(u, v)$ by (b).
Hence $c_H(u, v) = c_K(u, v)$, and all these inequalities are equalities.

(d) Let $S$ be a minimum $u$--$v$ cut of $K$ containing a $u$--$v$ path $\pi$.
By \cref{lem:zz-paths-in-q}(b), all vertices of $\pi$ lie in $Q$, so its edges are real edges.
Each of them has exactly one end in $S$, hence exactly one in $S \cap Q$, which is a minimum cut of $H$ by (c).
Conversely, let $B$ be a minimum cut of $H$ and $\pi$ a $u$--$v$ path of real edges in $\delta_H(B)$.
Each edge of $\pi$ has both ends in $Q$ and exactly one of them in $B = A(B) \cap Q$, so it lies in $\delta_K(A(B))$, and $A(B)$ is a minimum cut of $K$ by (c).

(e) By \cref{lem:zz-paths-in-q}(c), we have $d_K(x, y) \leq 2$.
If $xy \in E(K)$, it is a real edge.
Otherwise $x$ and $y$ have a common neighbor $z$.
If $z \in Q$, then $xz$ and $yz$ are real edges.
If $z \notin Q$, then $z$ lies in an ear $Y$ by \cref{lem:zz-ears}(c), and $xz$ and $yz$ are two distinct edges of $\delta_K(Y)$.
As $|\delta_K(Y)| = 2$, the attachments of $Y$ are $x$ and $y$, with $x \neq y$, and $H$ has the ear edge $xy$.
\end{proof}

\begin{lemma}[Uncrossing in $H$]
\label{lem:zz-uncrossing}
Let $S_1$ and $S_2$ be minimum $u$--$v$ cuts of $H$.
Then $S_1 \cap S_2$ and $S_1 \cup S_2$ are minimum $u$--$v$ cuts of $H$, and no edge of $H$ joins $S_1 \setminus S_2$ and $S_2 \setminus S_1$.
\end{lemma}

\begin{proof}
The proof of \cref{lem:uncrossing} counts one edge at a time, so it gives
\[
|\delta_H(S_1 \cap S_2)| + |\delta_H(S_1 \cup S_2)| = |\delta_H(S_1)| + |\delta_H(S_2)| - 2 j
\]
for the multigraph $H$, where $j$ is the number of edges, with multiplicity, between $S_1 \setminus S_2$ and $S_2 \setminus S_1$.
The sets $S_1 \cap S_2$ and $S_1 \cup S_2$ contain $u$ and not $v$, so each of them has at least $c_H(u, v)$ boundary edges.
The identity with $j \geq 0$ forces both values to equal $c_H(u, v)$ and $j = 0$.
\end{proof}

\begin{definition}[The ZZ instance of $(G, u, v)$]
\label{def:zz-instance}
For $x, y \in Q$ let $x \leq y$ if every minimum $u$--$v$ cut of $H$ that contains $y$ contains $x$.
Let $R$ be the set of real edges of $H$ whose ends are not tied.
The \emph{ZZ instance of} $(G, u, v)$ is $(Q, \leq, R, u, v)$.
\end{definition}

\begin{lemma}[Incomparable elements]
\label{lem:zz-incomparable}
The relation $\leq$ is a preorder on $Q$ with $u \leq x \leq v$ for every $x \in Q$, and $u < v$.
If $x, y \in Q$ are incomparable, no edge of $H$, real or ear, joins them.
Hence every edge of $R$ joins two strictly comparable elements, and $\{u, v\} \notin R$.
\end{lemma}

\begin{proof}
Reflexivity and transitivity follow from the definition.
Every minimum cut contains $u$, so $u \leq x$, and none contains $v$, so $x \leq v$.
A minimum $u$--$v$ cut of $H$ exists and contains $u$ but not $v$, so $v \leq u$ fails and $u < v$.
Let $x \parallel y$.
As $x \leq y$ fails, some minimum cut $S_2$ contains $y$ and not $x$.
As $y \leq x$ fails, some minimum cut $S_1$ contains $x$ and not $y$.
By \cref{lem:zz-uncrossing}, no edge joins $S_1 \setminus S_2 \ni x$ and $S_2 \setminus S_1 \ni y$.
An edge of $R$ joins ends that are not tied and, by what was just shown, comparable, hence strictly comparable.
We have $uv \notin E(K)$ since $d_K(u, v) = 2$, so $\{u, v\}$ is not a real edge.
\end{proof}

\begin{lemma}[Minimum cuts are down-sets]
\label{lem:zz-down-sets}
A set $D \subseteq Q$ is a minimum $u$--$v$ cut of $H$ if and only if $D$ is a down-set of $(Q, \leq)$ with $u \in D$ and $v \notin D$.
\end{lemma}

\begin{proof}
Let $D$ be a minimum cut, $y \in D$ and $x \leq y$.
By the definition of $\leq$, we have $x \in D$.
So $D$ is a down-set, and $u \in D$, $v \notin D$.
Conversely, let $D$ be a down-set with $u \in D$ and $v \notin D$, and let $y \in D$.
Some minimum cut contains $y$, since otherwise $v \leq y$ holds vacuously and $v \in D$.
The intersection $S_y$ of the minimum cuts containing $y$ is a minimum cut by repeated use of \cref{lem:zz-uncrossing}.
By the definition of $\leq$, it equals $\{x : x \leq y\}$, which lies in $D$ as $D$ is a down-set.
So $D = \bigcup_{y \in D} S_y$ is a union of minimum cuts, non-empty as $u \in D$, hence a minimum cut by \cref{lem:zz-uncrossing}.
\end{proof}

\begin{lemma}[P2]
\label{lem:zz-p2}
The ZZ instance $(Q, \leq, R, u, v)$ satisfies P2.
\end{lemma}

\begin{proof}
Let $x \parallel y$, so $x \neq y$.
By \cref{lem:zz-multigraph}(e), the elements $x$ and $y$ are joined by a real or an ear edge, or have a common neighbor $z \in Q$ with real edges $xz$ and $yz$.
The first two cases are excluded by \cref{lem:zz-incomparable}.
By the same lemma, $z$ is comparable with $x$ and with $y$.
If $x \leq z \leq y$ or $y \leq z \leq x$, then $x$ and $y$ would be comparable.
Hence $z \leq x$ and $z \leq y$, or $z \geq x$ and $z \geq y$.
If $z \sim x$, comparability of $z$ and $y$ would make $x$ and $y$ comparable, so $z$ is not tied with $x$, and likewise not with $y$.
Hence $z < x$ and $z < y$, or $z > x$ and $z > y$, and $xz$ and $yz$ join elements that are not tied, so they belong to $R$.
So $z$ is a lower or an upper witness.
\end{proof}

\begin{theorem}[Reduction]
\label{thm:zz-reduction}
Let $(G, u, v)$ be an instance of the open case with $n$ vertices.
Its ZZ instance $(Q, \leq, R, u, v)$ satisfies P2 and $u < v$, has $|Q| \leq n$ and $|R| \leq |E|$, and is computable in time polynomial in $n$.
We have $\cp(u, v) = c(u, v)$ if and only if it is a yes-instance of ZZ.
\end{theorem}

\begin{proof}
P2 is \cref{lem:zz-p2}, and $u < v$ and the form of $R$ are \cref{lem:zz-incomparable}.
We have $Q \subseteq V(K) \subseteq V$ and $R \subseteq E(K[Q]) \subseteq E$.

\emph{Answer.}
By \cref{prop:open-pairs}, $\cp(u, v) = c(u, v)$ if and only if some minimum $u$--$v$ cut of $G$ contains a $u$--$v$ path.
By \cref{lem:zz-bridges}(c) and \cref{lem:minimum-cuts}, this holds if and only if some minimum $u$--$v$ cut $S$ of $K$ contains a $u$--$v$ path.
By \cref{lem:zz-multigraph}(d), this holds if and only if some minimum $u$--$v$ cut $B$ of $H$ has a $u$--$v$ path of real edges inside $\delta_H(B)$.
By \cref{lem:zz-down-sets}, these sets $B$ are exactly the down-sets $D$ with $u \in D$ and $v \notin D$.
For such $D$, a real edge with exactly one end in $D$ joins elements that are not tied, because a down-set contains, with an element, its whole class, so the edge lies in $R$.
Conversely, an edge of $R$ with exactly one end in $D$ is a real edge in $\delta_H(D)$.
So the real edges inside $\delta_H(D)$ are exactly the edges of $G_D$, and the condition says that $u$ and $v$ are connected in $G_D$ for some such $D$.

\emph{Time.}
The bridges and $K$ take $O(n + |E|)$ time.
The set $Q$ takes one test $c_K(u, x) \geq 3$ per vertex, each by at most three augmenting-path searches.
The ears are the components of $K - Q$ by \cref{lem:zz-ears}(c), with their attachments read off their boundary edges.
This gives $H$, with $|Q| \leq n$ vertices and at most $|E| + n$ edges.
The value $c_K(u, v)$ takes one maximum-flow computation.
For the preorder, $x \leq y$ holds if $x = y$, $x = u$ or $y = v$.
Otherwise it holds if and only if no minimum $u$--$v$ cut of $H$ contains $y$ but not $x$, that is, if and only if the least $|\delta_H(S)|$ over all $S$ with $\{u, y\} \subseteq S \subseteq Q \setminus \{x, v\}$ exceeds $c_K(u, v)$.
This is one maximum-flow computation between $\{u, y\}$ and $\{x, v\}$.
So at most $|Q|^2$ flow computations in $H$ are needed, and $R$ follows.
All of this is polynomial in $n$.
\end{proof}

\subsection{The Switch Form of a ZZ Instance}
\label{sec:zz-construction}

\begin{definition}[Switch form of a ZZ instance]
\label{con:zz-switch-form}
Let $(Q, \leq, R, u, v)$ be a ZZ instance with $u < v$, and let $[u]$ and $[v]$ be the classes of $u$ and $v$, which are disjoint as $u < v$.
The \emph{switches} are the elements of $M = Q \setminus ([u] \cup [v])$, with the restriction of $\leq$, and $m = |M|$.
A real edge $ab$ with $a < b$ is of one of four kinds, namely $a \in [u]$, $b \in [v]$, or $a \in [u]$, $b \in M$, or $a \in M$, $b \in [v]$, or $a, b \in M$.
Indeed, its ends are strictly comparable, so they are not both in $[u]$ or both in $[v]$, and $a \in [v]$ would give $v \leq a < b \leq v$, and $b \in [u]$ would give $u \leq a < b \leq u$.
The \emph{nodes} are the vertex sets of the components of the graph on $[u] \cup [v]$ whose edges are the real edges of the first kind, and one new node $\rho(a, b)$ for every real edge $ab$ of the fourth kind.
We write $\nodeof(x)$ for the component of $x \in [u] \cup [v]$.
The terminals are $s = \nodeof(u)$ and $t = \nodeof(v)$.
For $w \in M$, we put
\begin{align*}
E_w &= \{\nodeof(a) : aw \in R,\ a \in [u]\} \cup \{\rho(a, w) : aw \in R,\ a \in M,\ a < w\},\\
F_w &= \{\nodeof(b) : wb \in R,\ b \in [v]\} \cup \{\rho(w, b) : wb \in R,\ b \in M,\ w < b\}.
\end{align*}
The middle edges are the pairs $(a, b)$ for the real edges $ab$ of the fourth kind, with private node $\rho(a, b)$.
\end{definition}

\begin{lemma}[Well defined, P2, sizes]
\label{lem:zz-switch-form}
\Cref{con:zz-switch-form} gives an instance of the general switch form, that is, every middle edge $(a, b)$ has $a < b$, and $\rho(a, b)$ lies in $F_a$ and $E_b$ and in no other side.
If the ZZ instance satisfies P2, so does its switch form.
The sizes satisfy $m \leq |Q| - 2$, $|\Omega| \leq |Q| + |R|$ and $\sum_{w \in M} (|E_w| + |F_w|) \leq 2|R|$, and the construction takes polynomial time.
\end{lemma}

\begin{proof}
The node $\rho(a, b)$ is put into $F_a$, by the second part of $F_w$ for $w = a$, and into $E_b$, by the second part of $E_w$ for $w = b$, and nowhere else.

For P2, let $x, y \in M$ be incomparable, in $M$ or equivalently in $Q$, and let $z$ be a common $R$-neighbor as in \cref{def:zz-zz}.
If $z < x$ and $z < y$, then $z \notin [v]$, since otherwise $v \leq z < x \leq v$.
For $z \in [u]$, the node $\nodeof(z)$ lies in $E_x \cap E_y$.
For $z \in M$, the pairs $(z, x)$ and $(z, y)$ are middle edges and $z$ is a lower middle witness.
If $z > x$ and $z > y$, then dually $z \notin [u]$.
For $z \in [v]$, the node $\nodeof(z)$ lies in $F_x \cap F_y$.
For $z \in M$, the element $z$ is an upper middle witness.

For the sizes, $[u]$ and $[v]$ are non-empty and disjoint, so $m \leq |Q| - 2$.
There are at most $|[u] \cup [v]|$ components and at most $|R|$ private nodes.
A real edge of the second or third kind puts one node into one side, one of the fourth kind puts its private node into two sides, and one of the first kind puts no node into a side.
\end{proof}

\begin{lemma}[Answers agree]
\label{lem:zz-answers}
A ZZ instance with $u < v$ is a yes-instance if and only if its switch form is a yes-instance.
\end{lemma}

\begin{proof}
(i) \emph{Down-sets.}
Let $D'$ range over the down-sets of $(Q, \leq)$ with $u \in D'$ and $v \notin D'$, and $D$ over the mode vectors.
We show that the maps $D' \mapsto D' \cap M$ and $D \mapsto [u] \cup D$ are inverse bijections.
The set $D'$ contains $[u]$, because $b \leq u \in D'$ for $b \in [u]$.
It misses $[v]$, because $x \in D' \cap [v]$ gives $v \leq x$ and $v \in D'$.
The set $D' \cap M$ is a down-set of $M$.
Conversely, let $D$ be a mode vector.
An element $x$ below some $b \in [u]$ lies in $[u]$, because $x \leq b \leq u \leq x$.
An element $x$ below some $d \in D$ lies in $[u]$, or in $M$ and then in $D$, but not in $[v]$, because $v \leq x \leq d$ would put $d$ into $[v]$.
So $[u] \cup D$ is a down-set, and it contains $u$ and not $v$.

(ii) \emph{Edges.}
Let $D' = [u] \cup D$, and let $\pk(w)$ be the side picked by $D$ for $w \in M$.
The edges of $G_{D'}$ are the following.
Every real edge of the first kind is one.
Every real edge $aw$ of the second kind with $w \notin D$ is one, and then $\nodeof(a) \in E_w = \pk(w)$.
Every real edge $wb$ of the third kind with $w \in D$ is one, and then $\nodeof(b) \in F_w = \pk(w)$.
Every real edge $zw$ of the fourth kind, with $z < w$, is one if $z \in D$ and $w \notin D$, because it needs exactly one end in $D$ and $w \in D$ forces $z \in D$, and then $\rho(z, w) \in F_z \cap E_w = \pk(z) \cap \pk(w)$.
In particular, an edge of $G_{D'}$ between $w \in M$ and $e \in [u] \cup [v]$ puts $\nodeof(e)$ into $\pk(w)$.

(iii) \emph{If $u$ and $v$ are connected in $G_{D'}$, then $D$ connects $s$ and $t$.}
Take a $u$--$v$ path of $G_{D'}$.
A maximal run of consecutive vertices in $[u] \cup [v]$ is joined by edges of the first kind, so all its vertices have the same node.
If the path has no vertex in $M$, then $\nodeof(u) = \nodeof(v)$, that is, $s = t$.
Otherwise let $w_1, \dots, w_h$ be its vertices in $M$ in order.
If $w_i$ and $w_{i+1}$ are adjacent on the path, then $\pk(w_i) \cap \pk(w_{i+1})$ contains the private node of their edge.
Otherwise the vertices between them form a run with one node $o$, and by (ii) the edges from $w_i$ and $w_{i+1}$ into that run put $o$ into $\pk(w_i)$ and $\pk(w_{i+1})$.
The run before $w_1$ contains $u$, so by (ii) we have $s = \nodeof(u) \in \pk(w_1)$, and likewise $t \in \pk(w_h)$.
So $\pk(w_1), \dots, \pk(w_h)$ connect $s$ and $t$.

(iv) \emph{If $D$ connects $s$ and $t$, then $u$ and $v$ are connected in $G_{D'}$.}
If $s = t$, then $u$ and $v$ lie in one component of the graph of the real edges of the first kind, all of which are edges of $G_{D'}$.
Otherwise let $X_1, \dots, X_h$ be picked sides, $X_i = \pk(w_i)$, with $s \in X_1$, $t \in X_h$ and $X_i \cap X_{i+1} \neq \emptyset$.
Suppose that a node $o$ lies in $\pk(w) \cap \pk(w')$ with $w \neq w'$.
Then $w$ and $w'$ are connected in $G_{D'}$, as we now show.
If $o$ is a component node, then by the construction and (ii), $w$ and $w'$ have $G_{D'}$-neighbors $e$ and $e'$ in the component $o$, and $e$ and $e'$ are joined by edges of the first kind.
If $o = \rho(z, w'')$, it lies only in $F_z$ and $E_{w''}$, so $\{w, w'\} = \{z, w''\}$, $z \in D$, $w'' \notin D$, and $zw''$ is an edge of $G_{D'}$.
Finally, $s \in \pk(w_1)$ is a component node, so $w_1$ has a $G_{D'}$-neighbor in the component $s \ni u$, which is joined to $u$ by edges of the first kind.
Likewise $w_h$ is connected to $v$.

By (i), (iii) and (iv), some $G_{D'}$ connects $u$ and $v$ if and only if some mode vector connects $s$ and $t$.
\end{proof}

\begin{lemma}[Valid paths and yes-instances]
\label{lem:zz-valid}
Consider an instance of the general switch form.
(a) Let $P$ be a valid path.
Then $U_P \cap L_P = \emptyset$, and a mode vector $D$ picks every member of $P$ if and only if $L_P \subseteq D \subseteq M \setminus U_P$.
In particular, the mode vector $L_P$ picks every member of $P$, and every mode vector that picks every member of $P$ connects $s$ and $t$.
(b) If a mode vector picks the sides of a sequence $Y_1, \dots, Y_h$ with $h \geq 1$, $s \in Y_1$, $t \in Y_h$ and $Y_i \cap Y_{i+1} \neq \emptyset$ for $i < h$, then a valid path with at most $h$ members exists.
In particular, if $s \neq t$ and some mode vector connects $s$ and $t$, then a valid path exists.
(c) The instance is a yes-instance if and only if $s = t$ or a valid path exists.
If $s = t$, every mode vector connects $s$ and $t$.
\end{lemma}

\begin{proof}
(a) If $x \in U_P \cap L_P$, then $w \leq x \leq z$ for some $w \in W_E$ and $z \in W_F$.
We have $w \neq z$ because the members belong to distinct switches, and $\{E_w, F_z\}$ is a conflict pair of two members, against validity.
The mode vector $D$ picks the F-member $F_w$ if and only if $w \in D$, so it picks all F-members if and only if $W_F \subseteq D$, that is, $L_P = \dn W_F \subseteq D$, as $D$ is a down-set.
The mode vector $D$ picks the E-member $E_w$ if and only if $w \notin D$, so it picks all E-members if and only if $W_E \cap D = \emptyset$, that is, $U_P \cap D = \emptyset$, because $x \in U_P \cap D$ gives $w \leq x$ for some $w \in W_E$, and then $w \in D$.
The set $L_P$ is a down-set disjoint from $U_P$, so it picks every member.
If $D$ picks every member, the members $X_1, \dots, X_q$ are picked sides with $s \in X_1$, $t \in X_q$ and $X_i \cap X_{i+1} \neq \emptyset$.

(b) Among all sequences of sides picked by this mode vector $D$ that start with a side containing $s$, end with a side containing $t$, and have consecutive sides meeting, let $X_1, \dots, X_{h'}$ be one with $h'$ least, so $h' \leq h$.
The sides are pairwise distinct, because $X_i = X_j$ with $i < j$ would allow us to drop $X_{i+1}, \dots, X_j$.
Hence they belong to distinct switches, since $D$ picks one side per switch.
Two picked sides $E_w$ and $F_z$ never form a conflict pair, because $w \notin D$, $z \in D$, and $w \leq z$ would put $w$ into $D$.
So $X_1, \dots, X_{h'}$ is a valid path.
If $s \neq t$ and $D$ connects $s$ and $t$, then such a sequence of picked sides exists by \cref{def:zz-valid}.

(c) If $s = t$, every mode vector connects $s$ and $t$ by \cref{def:zz-valid}.
If $s \neq t$, use (a) and (b).
\end{proof}
